\chapter{PageRank and Markov Chains on Graphs}
\label{ch:pagerank}

A web crawler in 1997 faced a brutal problem: it could fetch billions of pages, but it had no idea which ones \emph{mattered}. Keyword matching said the words ``Stanford University'' appeared on a hundred thousand pages -- which one should a searcher see first? The answer that reorganized an industry was simple to state. Treat the web as a directed graph, let a bored surfer click links at random forever, and ask: where does the surfer spend their time? Pages the surfer visits often are important. That long-run fraction of time is a \emph{stationary distribution} of a Markov chain, and computing it on a graph with billions of nodes is what this chapter is about.

We will build the random-surfer model from scratch, show that PageRank is the stationary distribution of a carefully patched Markov chain, prove that power iteration converges and pin down its rate to the spectral gap, develop personalized PageRank, and study how the algorithm bends (and resists bending) under adversarial link spam. Throughout, the graph view and the linear-algebra view are two readings of the same object, and we keep both open at once.

\begin{connectionbox}
This chapter sits on top of three others. The Markov-chain machinery -- stationary distributions, mixing, the Perron--Frobenius theorem -- is developed in the probability companion; we recall what we need. The matrix $\mathbf{M} = \Adjm^\top \mathbf{D}^{-1}$ and its eigenstructure connect to the spectral chapters (Ch.~31--34) and to the random-walk/electrical-network material (Ch.~35). Personalized PageRank is, almost literally, a discounted random walk, so the walk-and-resistance lens carries over directly.
\end{connectionbox}

\section{The web as a graph}
\label{ch36:web-graph}

Model the web as a directed graph $\Gph = (V, E)$. Each page is a vertex; a hyperlink from page $u$ to page $v$ is a directed edge $u \to v$. Let $n = \abs{V}$. For a vertex $u$ let $\deg^+(u)$ be its \emph{out-degree}, the number of links leaving $u$.

\begin{figure}[H]
\centering
\begin{tikzpicture}[->,>=Stealth,thick,
  vx/.style={circle,draw=nodeblue,fill=nodeblue!12,minimum size=8.5mm,font=\small\bfseries},
  vxTele/.style={circle,draw=structurepurple,fill=structurepurple!12,minimum size=8.5mm,font=\small\bfseries},
  link/.style={draw=nodeblue,thick},
  tele/.style={draw=structurepurple!70,dashed,thick}]

  % ================= Panel A: Link Graph =================
  \begin{scope}[shift={(0,0)}]
    \draw[rounded corners=6pt,fill=sourceblue!4,draw=sourceblue!25,thick] (-1.2,-2.7) rectangle (3.7,1.2);
    \node[font=\footnotesize\scshape,text=sourceblue] at (1.25,0.9) {Link Transitions $\mathbf{P}_{vu} = 1/\deg^+(u)$};

    \node[vx] (A) at (0,0) {$A$};
    \node[vx] (B) at (2.5,0) {$B$};
    \node[vx] (C) at (0,-1.8) {$C$};
    \node[vx] (D) at (2.5,-1.8) {$D$};

    \draw[link] (A) to[bend left=15] node[above=1pt,font=\tiny]{$\tfrac12$} (B);
    \draw[link] (B) to[bend left=15] node[below=1pt,font=\tiny]{$\tfrac12$} (A);
    \draw[link] (A) to node[left=1pt,font=\tiny]{$\tfrac12$} (C);
    \draw[link] (B) to node[right=1pt,font=\tiny]{$\tfrac12$} (D);
    \draw[link] (C) to[bend left=15] node[above=1pt,font=\tiny]{$1$} (D);
    \draw[link] (D) to[bend left=15] node[below=1pt,font=\tiny]{$\tfrac12$} (C);
    \draw[link] (D) to node[below left=1pt,font=\tiny]{$\tfrac12$} (B);
  \end{scope}

  % ================= Panel B: With Teleportation =================
  \begin{scope}[shift={(6.2,0)}]
    \draw[rounded corners=6pt,fill=structurepurple!4,draw=structurepurple!25,thick] (-1.2,-2.7) rectangle (3.7,1.2);
    \node[font=\footnotesize\scshape,text=structurepurple] at (1.25,0.9) {Google Matrix $\mathbf{G} = \alpha \mathbf{P} + (1-\alpha) v \mathbf{1}^\top$};

    \node[vxTele] (tA) at (0,0) {$A$};
    \node[vxTele] (tB) at (2.5,0) {$B$};
    \node[vxTele] (tC) at (0,-1.8) {$C$};
    \node[vxTele] (tD) at (2.5,-1.8) {$D$};

    % Standard links with damping alpha
    \draw[link] (tA) to[bend left=12] (tB);
    \draw[link] (tB) to[bend left=12] (tA);
    \draw[link] (tA) to (tC);
    \draw[link] (tB) to (tD);
    \draw[link] (tC) to[bend left=12] (tD);
    \draw[link] (tD) to[bend left=12] (tC);
    \draw[link] (tD) to (tB);

    % Central teleport distribution hub
    \node[rectangle,draw=structurepurple,fill=white,rounded corners=2pt,inner sep=2pt,font=\tiny\bfseries,text=structurepurple]
      (hub) at (1.25,-0.9) {Teleport $\tfrac{1-\alpha}{4}$};

    \draw[tele,<-] (tA) -- (hub);
    \draw[tele,<-] (tB) -- (hub);
    \draw[tele,<-] (tC) -- (hub);
    \draw[tele,<-] (tD) -- (hub);
  \end{scope}

\end{tikzpicture}
\caption{The four-page web and its PageRank Markov chain. Left: The raw link transition graph $\mathbf{P}$.
Surfers follow links with uniform probabilities $1/\deg^+(u)$. Right: The regularized Google matrix $\mathbf{G}$.
With probability $\alpha$, the surfer follows a link; with probability $1-\alpha$, they teleport to a uniformly
chosen page with probability $(1-\alpha)/n$. This guarantees that all transition probabilities are strictly
positive, yielding an irreducible, aperiodic Markov chain with a unique stationary distribution.}
\label{ch36:fig-web4}
\end{figure}

\begin{intuitionbox}
Forget content entirely. A link $u \to v$ is a tiny vote: ``$u$ thinks $v$ is worth visiting.'' But not all votes are equal. A link from a page that is itself widely linked-to should count for more than a link from an obscure page. This is recursive -- importance is defined in terms of importance -- and the random surfer is the trick that makes the recursion well-posed and computable.
\end{intuitionbox}

The recursive intuition has a long pedigree. Counting weighted citations to rank academic papers goes back to the 1970s (Pinski and Narin's \emph{influence weights}), and eigenvector centrality predates the web by decades. What Page and Brin added was a model -- the surfer -- that makes the eigenvector both interpretable and numerically tame on a graph the size of the web.

\section{The random-surfer model}
\label{ch36:surfer}

Picture a surfer standing on page $u$. With each tick of the clock they pick one of the $\deg^+(u)$ outgoing links uniformly at random and follow it. This is a random walk on the directed graph. Where the surfer is after $t$ steps is a random variable; the \emph{distribution} of that random variable is a column vector $x_t \in \R^n$ with $x_t(v) = \Pr[\text{surfer at } v \text{ at time } t]$.

\subsection{The transition matrix}

Define the column-stochastic \emph{transition matrix} $\mathbf{P} \in \R^{n\times n}$ by
\[
  \mathbf{P}_{vu} \;=\; \Pr[u \to v \text{ in one step}] \;=\;
  \begin{cases} \dfrac{1}{\deg^+(u)} & \text{if } (u,v) \in E,\\[1ex] 0 & \text{otherwise.}\end{cases}
\]
Column $u$ of $\mathbf{P}$ holds the surfer's next-step probabilities given they are at $u$; each nonzero column sums to $1$. Equivalently, with $\Adjm$ the directed adjacency matrix ($\Adjm_{uv}=1$ iff $u\to v$) and $\mathbf{D}=\operatorname{diag}(\deg^+(u))$,
\[
  \mathbf{P} \;=\; \Adjm^\top \mathbf{D}^{-1}.
\]
If we store distributions as \emph{column} vectors $x_t$, one step advances them by $x_{t+1} = \mathbf{P}\, x_t$. (Some texts use row vectors and a row-stochastic $\mathbf{P}^\top$; we fix the column convention and stick to it.)

\begin{example}[Transition matrix of Figure~\ref{ch36:fig-web4}]
\label{ch36:ex-Pmatrix}
Order the vertices $A,B,C,D$. Out-edges: $A\to\{B,C\}$, $B\to\{A,D\}$, $C\to\{D\}$, $D\to\{B,C\}$. So
\[
  \mathbf{P} =
  \begin{pmatrix}
    0 & \tfrac12 & 0 & 0\\
    \tfrac12 & 0 & 0 & \tfrac12\\
    \tfrac12 & 0 & 0 & \tfrac12\\
    0 & \tfrac12 & 1 & 0
  \end{pmatrix},
  \qquad
  \begin{array}{l}
  \text{column } A:\ \text{mass splits to } B,C\\
  \text{column } B:\ \text{mass splits to } A,D\\
  \text{column } C:\ \text{all mass to } D\\
  \text{column } D:\ \text{mass splits to } B,C
  \end{array}
\]
Each column sums to $1$, so $\mathbf{P}$ is column-stochastic.
\end{example}

\subsection{Two diseases: dangling nodes and traps}

The raw walk has two pathologies that break the ``long-run fraction of time'' story.

\begin{trapbox}
\textbf{Dangling nodes (rank sinks).} A page with no outgoing links ($\deg^+(u)=0$) has an all-zero column in $\mathbf{P}$. The surfer who arrives there \emph{vanishes} -- probability leaks out of the system, and $\mathbf{P}$ is no longer stochastic. PDFs, image pages, and pages the crawler has not yet expanded are all dangling.

\textbf{Traps (disconnected cliques).} Even with every out-degree positive, a group of pages that link only among themselves swallows all probability that ever enters it. In Markov language the chain is \emph{reducible} or \emph{periodic}, and a unique stationary distribution need not exist.
\end{trapbox}

Both diseases have the same cure, and it is the surfer's one extra habit: \textbf{teleportation}.

\section{PageRank as a stationary distribution}
\label{ch36:stationary}

Give the surfer a coin with bias $\alpha$ (the \emph{damping factor}, classically $\alpha=0.85$). At each step:
\begin{itemize}
  \item with probability $\alpha$, follow a uniformly random out-link (the old behavior);
  \item with probability $1-\alpha$, \emph{teleport} to a page chosen from a distribution $v$ (uniform, $v=\tfrac1n\mathbf{1}$, in the classic model).
\end{itemize}
A dangling node is handled by saying: if there are no out-links, teleport with probability $1$. Algebraically, patch each all-zero column of $\mathbf{P}$ to the teleport vector $v$, giving a genuinely column-stochastic $\mathbf{P}'$. The \emph{Google matrix} is then
\[
  \boxed{\;\mathbf{G} \;=\; \alpha\, \mathbf{P}' \;+\; (1-\alpha)\, v\,\mathbf{1}^\top.\;}
\]
Here $v\,\mathbf{1}^\top$ is the rank-one matrix every column of which equals $v$: from anywhere, teleport lands you in $v$. The PageRank vector $\pi$ is the stationary distribution of $\mathbf{G}$:
\[
  \boxed{\;\mathbf{G}\,\pi = \pi, \qquad \pi \ge 0, \qquad \mathbf{1}^\top \pi = 1.\;}
\]

\begin{intuitionbox}
Teleportation is a humility tax on the link structure. A fraction $1-\alpha$ of the time the surfer ignores links entirely and jumps to a random page. This single modification resolves both pathologies: it provides a smooth exit from dead-end dangling nodes, dissolves traps that would otherwise swallow probability inside closed cliques, and ensures that every entry of the Google matrix $\mathbf{G}$ is strictly positive. By the Perron--Frobenius theorem, this strict positivity guarantees a unique stationary distribution that power iteration is guaranteed to reach.
\end{intuitionbox}

\subsection{Why a unique stationary distribution exists}

\begin{theorem}[Existence and uniqueness of PageRank]
\label{ch36:thm-existence}
For any $0 < \alpha < 1$ and any probability vector $v > 0$ (all entries positive), the Google matrix $\mathbf{G}$ has a unique stationary distribution $\pi$ with $\pi > 0$, and it is the unique eigenvector of $\mathbf{G}$ with eigenvalue $1$ up to scaling.
\end{theorem}

\begin{whybox}
The proof is a clean application of Perron--Frobenius. Recall the statement: if $\mathbf{M}$ is a nonnegative \emph{primitive} matrix (some power is strictly positive), then its spectral radius $\rho$ is a simple eigenvalue, the only eigenvalue of modulus $\rho$, and its eigenvector can be chosen strictly positive. We check $\mathbf{G}$ qualifies.
\end{whybox}

\begin{proof}
Because $v > 0$ and $\alpha < 1$, every entry of $\mathbf{G} = \alpha\mathbf{P}' + (1-\alpha)v\mathbf{1}^\top$ satisfies $\mathbf{G}_{ij} \ge (1-\alpha)v_i > 0$. So $\mathbf{G}$ is strictly positive, hence irreducible and primitive (it is its own positive power). $\mathbf{G}$ is column-stochastic: $\mathbf{1}^\top \mathbf{G} = \alpha\,\mathbf{1}^\top\mathbf{P}' + (1-\alpha)(\mathbf{1}^\top v)\mathbf{1}^\top = \alpha\mathbf{1}^\top + (1-\alpha)\mathbf{1}^\top = \mathbf{1}^\top$, using $\mathbf{1}^\top \mathbf{P}' = \mathbf{1}^\top$ and $\mathbf{1}^\top v = 1$. A column-stochastic matrix has spectral radius exactly $1$, with $\mathbf{1}$ a left eigenvector for eigenvalue $1$. By Perron--Frobenius, $1$ is a simple eigenvalue, it is the unique eigenvalue of modulus $1$, and the corresponding right eigenvector $\pi$ can be chosen strictly positive. Normalize $\mathbf{1}^\top\pi = 1$. Uniqueness of the normalized positive eigenvector follows from simplicity.
\end{proof}

\begin{connectionbox}
This is the same Perron--Frobenius theorem that gives expander graphs (Ch.~34) their spectral gap and that governs the convergence of \emph{any} aperiodic irreducible Markov chain in the probability book. The strict positivity of $\mathbf{G}$ is the discrete-graph analogue of an ergodic chain: every state reaches every other in one step. PageRank is, then, exactly the long-run occupancy of an ergodic chain -- the ``fraction of time the surfer spends'' story made rigorous.
\end{connectionbox}

\subsection{Handling dangling nodes without densifying}

Patching every dangling column to $v$ literally would destroy sparsity. The fix is to keep $\mathbf{P}$ sparse and carry the dangling correction as a scalar. Let $D \subseteq V$ be the dangling set and let $d \in \{0,1\}^n$ be its indicator. Then $\mathbf{P}' = \mathbf{P} + v\, d^\top$ (mass landing in $D$ is redirected to $v$), and
\[
  \mathbf{G}\,x = \alpha\mathbf{P}x \;+\; \big(\alpha\, (d^\top x) + (1-\alpha)\big)\, v ,
  \qquad \text{whenever } \mathbf{1}^\top x = 1.
\]
The term $d^\top x = \sum_{u\in D} x_u$ is the total mass currently on dangling pages; it is a single dot product. We never form $\mathbf{G}$ or $\mathbf{P}'$ explicitly. This is the identity that makes PageRank scale.

\section{Power iteration and convergence}
\label{ch36:power}

How do we actually solve $\mathbf{G}\pi = \pi$ when $n$ is in the billions? We do \emph{not} invert anything. We iterate the surfer's step:
\[
  x^{(0)} = v, \qquad x^{(k+1)} = \mathbf{G}\, x^{(k)}.
\]
This is power iteration. Because $\mathbf{G}$ has a unique dominant eigenvalue $1$ with a spectral gap below it, $x^{(k)} \to \pi$.

\begin{algorithm}[H]
\caption{PageRank by power iteration (sparse, dangling-aware)}
\label{ch36:alg-pagerank}
\KwIn{Out-adjacency lists of $\Gph$; damping $\alpha$; teleport vector $v$; tolerance $\eps$}
\KwOut{PageRank vector $\pi$}
$x \gets v$\;
\Repeat{$\norm{x - x_{\mathrm{old}}}_1 \le \eps$}{
  $x_{\mathrm{old}} \gets x$\;
  $y \gets \mathbf{0}$ \tcp*{accumulate $\alpha\,\mathbf{P}x$}
  \ForEach{vertex $u$ with $\deg^+(u) > 0$}{
    \ForEach{$u \to w$}{ $y_w \gets y_w + \alpha\, x_u / \deg^+(u)$\;}
  }
  $s \gets \sum_{u \in D} x_u$ \tcp*{dangling mass}
  $\beta \gets \alpha\, s + (1-\alpha)$ \tcp*{redistributed via teleport}
  $x \gets y + \beta\, v$ \tcp*{$\mathbf{1}^\top x = 1$ automatically}
}
\Return $x$\;
\end{algorithm}

\begin{complexitybox}
Each iteration touches every edge once and every vertex a constant number of times: $\Th(n + m)$ time, where $m = \abs{E}$, and $\Th(n)$ extra space for $x, y$. The number of iterations to reach $\ell_1$-error $\eps$ is $\Oh\!\big(\log(1/\eps) / \log(1/\alpha)\big)$ (proved below), \emph{independent of $n$}. For $\alpha = 0.85$, each iteration cuts the error by $0.85$, so $\approx 50$--$100$ iterations give roughly $4$--$7$ correct digits on any graph (machine precision near $10^{-16}$ takes a couple hundred). Total: $\Oh\!\big((n+m)\,\log(1/\eps)/\log(1/\alpha)\big)$.
\end{complexitybox}

\subsection{The convergence rate is the damping factor}

The headline fact is that the second-largest eigenvalue of $\mathbf{G}$ is bounded by $\alpha$ \emph{regardless of the link graph}. The teleport term manufactures a spectral gap of at least $1-\alpha$.

\begin{theorem}[Subdominant eigenvalue]
\label{ch36:thm-eigen}
Let $\lambda_1 = 1, \lambda_2, \dots, \lambda_n$ be the eigenvalues of the column-stochastic matrix $\mathbf{P}'$, ordered so $\abs{\lambda_2}\ge\dots\ge\abs{\lambda_n}$. Then the eigenvalues of $\mathbf{G} = \alpha\mathbf{P}' + (1-\alpha)v\mathbf{1}^\top$ are $1$ together with $\alpha\lambda_2, \alpha\lambda_3, \dots, \alpha\lambda_n$. In particular the subdominant eigenvalue of $\mathbf{G}$ has modulus at most $\alpha$.
\end{theorem}

\begin{proof}
Let $\mathbf{P}'$ have right eigenvector $u_1 = \pi'$ (its own stationary vector, eigenvalue $1$, normalized $\mathbf{1}^\top u_1 = 1$). The left eigenvector of $\mathbf{P}'$ for eigenvalue $1$ is $\mathbf{1}^\top$. For any other eigenvalue $\lambda_i$ ($i \ge 2$) with right eigenvector $u_i$, left-eigenvector orthogonality gives $\mathbf{1}^\top u_i = 0$. Now compute $\mathbf{G}u_i = \alpha\mathbf{P}'u_i + (1-\alpha)v(\mathbf{1}^\top u_i) = \alpha\lambda_i u_i + 0 = \alpha\lambda_i\, u_i$. So each $u_i$ ($i\ge2$) is an eigenvector of $\mathbf{G}$ with eigenvalue $\alpha\lambda_i$. And $\mathbf{G}\pi = \pi$ supplies the eigenvalue-$1$ direction. Since $\mathbf{P}'$ stochastic forces $\abs{\lambda_i}\le 1$, we get $\abs{\alpha\lambda_i}\le\alpha$.
\end{proof}

\begin{theorem}[Power-iteration convergence rate]
\label{ch36:thm-converge}
With $x^{(k)} = \mathbf{G}^k v$ and $\pi$ the PageRank vector,
\[
  \norm{x^{(k)} - \pi}_1 \;\le\; 2\,\alpha^{k}.
\]
Thus $\Oh(\log(1/\eps)/\log(1/\alpha))$ iterations suffice for $\ell_1$-error $\eps$. Figure~\ref{ch36:fig-converge} plots this decay for several damping factors.
\end{theorem}

\begin{proof}
Write $e^{(k)} = x^{(k)} - \pi$. Since $\mathbf{G}\pi = \pi$, $e^{(k+1)} = \mathbf{G}e^{(k)}$, so $e^{(k)} = \mathbf{G}^k e^{(0)}$. The error vector lives in the subspace $W = \{w : \mathbf{1}^\top w = 0\}$, because both $x^{(k)}$ and $\pi$ are probability vectors and so $\mathbf{1}^\top e^{(k)} = 0$. On $W$, the action of $\mathbf{G}$ is governed by the eigenvalues $\alpha\lambda_2,\dots$ of Theorem~\ref{ch36:thm-eigen}, all of modulus $\le\alpha$. The cleanest bound is the $\ell_1$ contraction $\norm{\mathbf{G}w}_1 \le \alpha\norm{w}_1$ for $w \in W$ (proved in the next Why-box). Hence $\norm{e^{(k)}}_1 \le \alpha^k \norm{e^{(0)}}_1 \le 2\alpha^k$, since $\norm{x^{(0)}-\pi}_1 \le \norm{x^{(0)}}_1 + \norm{\pi}_1 = 2$.
\end{proof}

\begin{whybox}
The $\ell_1$-contraction statement $\norm{\mathbf{G}w}_1 \le \alpha\norm{w}_1$ for zero-sum $w$ has a clean probabilistic reading. Couple two surfers with current distributions differing by $w$; with probability $1-\alpha$ both teleport to the \emph{same} place $v$ and their distributions agree there, killing that fraction of the disagreement; with probability $\alpha$ they take a link step, which (being stochastic) does not \emph{increase} total-variation distance. So one step shrinks the disagreement by at least the teleport fraction $1-\alpha$ -- equivalently, retains at most $\alpha$ of it. This coupling \emph{is} the spectral-gap bound, with no diagonalization, and it works even when $\mathbf{G}$ is not diagonalizable.
\end{whybox}

\begin{trapbox}
A bigger $\alpha$ (say $0.95$) makes PageRank track link structure more closely but \emph{slows convergence} ($0.95^k$ decays lazily) and makes $\mathbf{G}$ more ill-conditioned and more spam-sensitive. A smaller $\alpha$ ($0.7$) converges fast and resists spam but washes out the link signal toward uniform. The classic $0.85$ is an engineering compromise, not a theorem.
\end{trapbox}

\begin{figure}[H]
\centering
\begin{tikzpicture}
\begin{axis}[
  width=0.82\linewidth, height=5.4cm,
  title={Geometric Convergence of Power Iteration ($2\alpha^k$)},
  title style={font=\small},
  xlabel={iteration $k$}, ylabel={error bound $2\alpha^k$},
  ymode=log, log basis y=10,
  xmin=0, xmax=100, ymin=1e-7, ymax=3,
  grid=both, grid style={gray!18},
  legend pos=north east, legend cell align=left,
  tick label style={font=\footnotesize}, label style={font=\small},
]
  \addplot[sourceblue,very thick,domain=0:100,samples=60] {2*0.70^x};
  \addplot[targetgreen,very thick,domain=0:100,samples=60] {2*0.85^x};
  \addplot[trapred,very thick,domain=0:100,samples=60] {2*0.95^x};
  \legend{$\alpha=0.70$,$\alpha=0.85$,$\alpha=0.95$}
\end{axis}
\end{tikzpicture}
\caption{The guaranteed $\ell_1$-error bound $\norm{x^{(k)}-\pi}_1 \le 2\alpha^k$ of Theorem~\ref{ch36:thm-converge}, on a log scale. The slope is $\log\alpha$ and nothing about the graph enters: the teleport-manufactured gap $1-\alpha$ sets the decay. At $\alpha=0.85$ roughly four dozen iterations buy three digits; $\alpha=0.95$ needs about three times as many for the same accuracy.}
\label{ch36:fig-converge}
\end{figure}

\subsection{The same thing as a linear solve}
\label{ch36:linsolve}

Power iteration is not the only way in. Subtracting the teleport term turns the stationary equation into a sparse linear system. Starting from $\pi = \alpha\mathbf{P}'\pi + (1-\alpha)v$ and rearranging,
\[
  (\mathbf{I} - \alpha\mathbf{P}')\,\pi \;=\; (1-\alpha)\,v.
\]
The coefficient matrix $\mathbf{I}-\alpha\mathbf{P}'$ is \emph{strictly diagonally dominant by columns} (each off-diagonal column sum is $\alpha < 1$), hence invertible and well-conditioned away from $\alpha = 1$. So PageRank is the solution of an $n\times n$ sparse linear system, and any Krylov solver -- GMRES, BiCGSTAB -- or a stationary iteration like Gauss--Seidel applies. Power iteration is precisely the \emph{Jacobi} iteration for this system, which is why its convergence factor is the spectral radius $\alpha$ of the iteration matrix $\alpha\mathbf{P}'$.

\begin{whybox}
Why bother, if power iteration already converges in $\Oh(\log(1/\eps)/\log(1/\alpha))$ steps? Two reasons. First, when $\alpha$ is pushed near $1$ (sharper link sensitivity), the Jacobi factor $\alpha$ approaches $1$ and power iteration crawls; a Krylov method's convergence depends on the full spectrum of $\mathbf{I}-\alpha\mathbf{P}'$, not merely its spectral radius, and can leap ahead. Second, the linear-system framing lets you reuse decades of preconditioning and parallel-solver engineering. The catch is memory: Krylov methods store several length-$n$ vectors and lose the clean ``one matrix--vector product per step'' structure that makes power iteration straightforward to distribute across a cluster. At web scale that simplicity usually wins, which is why production PageRank stayed close to power iteration even though the linear-algebra menu is larger.
\end{whybox}

\begin{connectionbox}
The matrix $\mathbf{I}-\alpha\mathbf{P}'$ is the discrete resolvent of the walk, and on an undirected graph it is a reweighted, shifted Laplacian. The fast Laplacian solvers of the spectral chapters (near-linear-time solves via combinatorial preconditioners) therefore apply to undirected PageRank directly -- another bridge between the Markov view here and the spectral-graph machinery of Ch.~31--34. The convergence factor $\alpha$ is exactly $1$ minus the spectral gap, so a slow-mixing chain and an ill-conditioned solve are the same phenomenon wearing two hats.
\end{connectionbox}

\section{Tracing PageRank by hand}
\label{ch36:trace}

Let us run Algorithm~\ref{ch36:alg-pagerank} on Figure~\ref{ch36:fig-web4}, with $\alpha = 0.85$ and uniform teleport $v = \tfrac14(1,1,1,1)^\top$. There are no dangling nodes here, so $s = 0$ and $\beta = (1-\alpha) = 0.15$ at every step; each step is $x \gets 0.85\,\mathbf{P}x + 0.0375\,\mathbf{1}$ (since $\beta v = 0.15\cdot\tfrac14\mathbf{1} = 0.0375\,\mathbf{1}$).

Start at $x^{(0)} = (0.25,\,0.25,\,0.25,\,0.25)^\top$. Recall from Example~\ref{ch36:ex-Pmatrix}: $A$ feeds $B,C$; $B$ feeds $A,D$; $C$ feeds $D$; $D$ feeds $B,C$.

\textbf{Step 1.} Compute $\mathbf{P}x^{(0)}$ entrywise. Mass into each node:
\begin{align*}
(\mathbf{P}x)_A &= \tfrac12 x_B = 0.125, &
(\mathbf{P}x)_B &= \tfrac12 x_A + \tfrac12 x_D = 0.25,\\
(\mathbf{P}x)_C &= \tfrac12 x_A + \tfrac12 x_D = 0.25, &
(\mathbf{P}x)_D &= \tfrac12 x_B + x_C = 0.375.
\end{align*}
Then $x^{(1)} = 0.85\,(0.125, 0.25, 0.25, 0.375)^\top + 0.0375\,\mathbf{1}$. Computing:
\[
  x^{(1)} = (0.14375,\; 0.25,\; 0.25,\; 0.35625)^\top.
\]
(Check: the entries sum to $1.0$.)

\textbf{Step 2.} Now feed $x^{(1)}$ through $\mathbf{P}$:
\begin{align*}
(\mathbf{P}x^{(1)})_A &= \tfrac12(0.25) = 0.125, &
(\mathbf{P}x^{(1)})_B &= \tfrac12(0.14375) + \tfrac12(0.35625) = 0.25,\\
(\mathbf{P}x^{(1)})_C &= 0.25, &
(\mathbf{P}x^{(1)})_D &= \tfrac12(0.25) + 0.25 = 0.375.
\end{align*}
So $x^{(2)} = 0.85(0.125,0.25,0.25,0.375)^\top + 0.0375\,\mathbf{1} = (0.14375, 0.25, 0.25, 0.35625)^\top$.

We hit a fixed point already: $x^{(2)} = x^{(1)}$. (This graph is symmetric enough that the $\mathbf{P}x$ profile stabilizes in one step; in general it would take a few dozen.) The PageRank vector is
\[
  \pi \approx (0.1438,\; 0.2500,\; 0.2500,\; 0.3563)^\top.
\]

\begin{labbox}
Read the answer. $D$ ranks highest ($0.356$): it receives from both $B$ and $C$, and $C$ pours \emph{all} its mass into $D$. $B$ and $C$ tie at $0.25$. $A$ is lowest ($0.144$): its only inbound link is from $B$, and that link is shared ($B$ splits to $A$ and $D$). Notice importance flowed against your first instinct -- $A$ has the same out-degree as $B$ and $D$, but out-degree does not help you; \emph{in}-links from important pages do.

Now re-run it yourself with $\alpha = 0.5$. You should see the spread between $\pi_D$ and $\pi_A$ shrink: more teleportation pulls every score toward the uniform $0.25$. Then add a dangling node $E$ (a page $A$ links to but that links nowhere) and watch $s = x_E$ become nonzero, feeding $\beta$.
\end{labbox}

\begin{historybox}
Sergey Brin and Larry Page described PageRank in ``The Anatomy of a Large-Scale Hypertextual Web Search Engine'' (1998) and the companion tech report ``The PageRank Citation Ranking: Bringing Order to the Web.'' The name puns on Larry Page. The mathematics is older: it is power iteration on a stochastic matrix, a technique von Mises studied in the 1920s, applied to an idea (eigenvector centrality) that sociologists and bibliometricians had circled for decades. The genius was systems-level -- recognizing that the teleport-regularized chain both made the eigenvector unique \emph{and} kept power iteration cheap on a web-scale sparse graph, then building the crawl-and-index pipeline around it.
\end{historybox}

\section{Personalized PageRank}
\label{ch36:ppr}

The teleport vector $v$ is a free parameter we set to uniform. Set it to something else and the whole ranking tilts toward that something. \emph{Personalized PageRank} (PPR) takes $v = e_s$, the indicator of a single source node $s$ (or a small set): every teleport sends the surfer \emph{home} to $s$. The stationary distribution then measures importance \emph{as seen from $s$} -- nodes a short, well-connected walk away from $s$ score high.

\[
  \boxed{\;\mathrm{ppr}_s \;=\; (1-\alpha)\,\big(\mathbf{I} - \alpha\mathbf{P}'\big)^{-1} e_s.\;}
\]

\begin{whybox}
Where does the closed form come from? The stationary equation $x = \alpha\mathbf{P}'x + (1-\alpha)e_s$ rearranges to $(\mathbf{I}-\alpha\mathbf{P}')x = (1-\alpha)e_s$. Since $\alpha < 1$ and the column-stochastic $\mathbf{P}'$ has spectral radius $1$, the matrix $\mathbf{I}-\alpha\mathbf{P}'$ is invertible with Neumann series $(\mathbf{I}-\alpha\mathbf{P}')^{-1} = \sum_{k\ge0}\alpha^k\mathbf{P}'^{\,k}$. So
\[
  \mathrm{ppr}_s = (1-\alpha)\sum_{k=0}^{\infty}\alpha^k\,\mathbf{P}'^{\,k} e_s.
\]
Read the $k$-th term: $\mathbf{P}'^{\,k}e_s$ is the distribution after $k$ link-steps from $s$, weighted by $(1-\alpha)\alpha^k$ -- exactly the probability that a teleport-home surfer takes precisely $k$ steps before the next teleport (a geometric$(1-\alpha)$ walk length). PPR is the \emph{expected discounted visit count} of short random walks launched from $s$. This is the discounted-walk view that connects straight to the random-walk chapter.
\end{whybox}

\begin{connectionbox}
That Neumann series is the discrete cousin of the resolvent $(\mathbf{I}-\alpha\mathbf{P})^{-1}$ and, on undirected graphs, of the Laplacian Green's function from Ch.~35. PPR vectors are the building block of local clustering: the Andersen--Chung--Lang algorithm sweeps a PPR vector to find a low-conductance cut \emph{near} $s$ without touching the whole graph -- the local analogue of Cheeger's inequality (Ch.~33). PPR is also how ``related pages,'' friend recommendation in social graphs, and many graph-neural-network propagations (APPNP) are computed.
\end{connectionbox}

\subsection{Computing PPR locally: push}

For a single source we rarely want the whole vector; we want the few hundred high-PPR nodes near $s$. The Andersen--Chung--Lang \emph{push} algorithm maintains an estimate $p$ and a residual $r$ (undelivered mass), and repeatedly pushes residual out of whichever node holds too much, touching only the explored frontier.

\begin{algorithm}[H]
\caption{Approximate Personalized PageRank by push (Andersen--Chung--Lang)}
\label{ch36:alg-push}
\KwIn{source $s$; damping $\alpha$; threshold $\tau$; out-adjacency of $\Gph$}
\KwOut{estimate $p \approx \mathrm{ppr}_s$ with bounded residual}
$p \gets \mathbf{0}$; \quad $r \gets e_s$\;
\While{$\exists\, u$ with $r_u \ge \tau\,\deg^+(u)$}{
  $p_u \gets p_u + (1-\alpha)\, r_u$ \tcp*{settle the keep-fraction at $u$}
  $m \gets \alpha\, r_u$; \quad $r_u \gets 0$\;
  \ForEach{$u \to w$}{ $r_w \gets r_w + m / \deg^+(u)$ \tcp*{spread the walk-fraction}}
}
\Return $p$\;
\end{algorithm}

\begin{whybox}
\emph{Correctness invariant.} Throughout, $\mathrm{ppr}_s = p + \mathrm{ppr}_r$, where $\mathrm{ppr}_r = \sum_u r_u\,\mathrm{ppr}_u$ is the PPR of the leftover residual (linearity of PPR in its teleport vector). Each push moves a node's residual into the part already accounted for by $p$ plus residual sent downstream, preserving the invariant exactly. When all residuals are below threshold, $p$ approximates $\mathrm{ppr}_s$ with $\ell_\infty$ error proportional to $\tau$.

\emph{Complexity.} Each push at $u$ settles at least $(1-\alpha)\tau\deg^+(u)$ mass permanently into $p$, and total settled mass is $\le 1$. So the sum of $\deg^+(u)$ over all pushes is $\Oh\!\big(\tfrac{1}{(1-\alpha)\tau}\big)$ -- a bound \emph{independent of $n$}. The work touches only nodes within reach of $s$, which is why push powers local clustering on graphs too large to load.
\end{whybox}

\begin{applicationbox}
Personalized PageRank is the quiet workhorse behind a lot of products: Twitter's ``Who To Follow'' used PPR on the follow graph; Pinterest's Pixie serves related-pin recommendations from random-walk/PPR scores at interactive latency; e-commerce ``customers also viewed'' and graph-database proximity queries are PPR in disguise. The push algorithm's $n$-independent cost is exactly what makes these feasible on billion-edge graphs.
\end{applicationbox}

\subsection{Tracing push by hand}
\label{ch36:push-trace}

Run Algorithm~\ref{ch36:alg-push} on Figure~\ref{ch36:fig-web4} with source $s = A$, damping $\alpha = 0.85$, and a generous threshold $\tau = 0.05$. The push rule fires at $u$ whenever $r_u \ge \tau\deg^+(u)$; settling deposits $(1-\alpha)r_u = 0.15\,r_u$ into $p_u$ and spreads $\alpha r_u = 0.85\,r_u$ across $u$'s out-edges. Recall the out-edges: $A\to\{B,C\}$, $B\to\{A,D\}$, $C\to\{D\}$, $D\to\{B,C\}$, each splitting mass into two halves except $C$, which sends everything to $D$.

\begin{figure}[H]
\centering
\begin{tikzpicture}[->,>=Stealth,node distance=2.2cm,thick,
  vx/.style={circle,draw=nodeblue,fill=blue!6,minimum size=9mm,font=\small},
  hot/.style={circle,draw=trapred,fill=red!10,minimum size=9mm,font=\small}]
  \node[hot] (A) {$A$};
  \node[vx,right=of A] (B) {$B$};
  \node[vx,below=1.5cm of A] (C) {$C$};
  \node[vx,below=1.5cm of B] (D) {$D$};
  \draw (A) to[bend left=12] node[above,font=\scriptsize,black]{$.425$} (B);
  \draw (A) -- node[left,font=\scriptsize,black]{$.425$} (C);
  \draw[gray!50] (B) to[bend left=12] (A);
  \draw[gray!50] (B) -- (D);
  \draw[gray!50] (C) to[bend left=12] (D);
  \draw[gray!50] (D) to[bend left=12] (C);
  \draw[gray!50] (D) -- (B);
\end{tikzpicture}
\caption{The first push, out of the source $A$. Its residual $r_A = 1$ settles $0.15$ into $p_A$ and sends $\alpha\cdot\tfrac12 = 0.425$ down each of $A$'s two out-edges. Gray edges are not yet active; the frontier grows outward from $s$.}
\label{ch36:fig-push}
\end{figure}

\textbf{Initialize.} $p = \mathbf{0}$, $r = e_A = (1,0,0,0)$. Only $A$ exceeds threshold ($r_A = 1 \ge 0.05\cdot 2$).

\textbf{Push $A$.} Settle $p_A \mathrel{+}= 0.15\cdot 1 = 0.15$; the walk-fraction is $m = 0.85$, split as $0.425$ to each of $B, C$; zero out $r_A$. Now
\[
  p = (0.15,\,0,\,0,\,0), \qquad r = (0,\,0.425,\,0.425,\,0).
\]

\textbf{Push $B$} ($r_B = 0.425 \ge 0.1$). Settle $p_B \mathrel{+}= 0.15\cdot 0.425 = 0.06375$; spread $m = 0.85\cdot 0.425 = 0.36125$ as $0.180625$ to each of $A, D$:
\[
  p = (0.15,\,0.06375,\,0,\,0), \quad r = (0.180625,\,0,\,0.425,\,0.180625).
\]

\textbf{Push $C$} ($r_C = 0.425 \ge 0.05$, since $\deg^+(C)=1$). Settle $p_C \mathrel{+}= 0.06375$; all of $m = 0.36125$ goes to $D$ (its only out-edge):
\[
  p = (0.15,\,0.06375,\,0.06375,\,0), \quad r = (0.180625,\,0,\,0,\,0.54125).
\]

After three pushes $p$ already orders the near neighborhood of $A$: the source keeps the most settled mass, $B$ and $C$ tie (each one hop from $A$), and $D$ is filling up in the residual as the two-hop sink. Continuing to push $D$ and the leftover residual at $A$ drives every residual under $\tau\deg^+$, at which point $p$ is the local PPR estimate. Notice we never touched a node outside $\{A,B,C,D\}$ -- on a billion-node web the same run would still settle only a handful of vertices.

\begin{labbox}
Carry the trace to completion. Keep pushing the largest over-threshold residual until none remain, and record the final $p$; normalize it and compare against $\mathrm{ppr}_A$ obtained from the closed form $(1-\alpha)(\mathbf{I}-\alpha\mathbf{P}')^{-1}e_A$ (here $\mathbf{P}'=\mathbf{P}$ since the graph has no dangling node; use any linear-algebra library). Then lower $\tau$ to $0.005$ and watch the number of pushes grow roughly like $1/\tau$, exactly as the complexity bound predicts. Finally, re-seed with $s = D$ and confirm the estimate now concentrates on $B, C$ -- the nodes $D$ feeds -- showing how PPR is a \emph{directional}, source-anchored centrality.
\end{labbox}

\section{Spam and robustness}
\label{ch36:spam}

PageRank's openness -- anyone can create pages and links -- is also its attack surface. Because a page's score is fed by its in-links, an adversary who controls many pages can sculpt link structure to inflate a target's rank. This is \emph{link spam}.

\subsection{The link farm}

\begin{figure}[H]
\centering
\begin{tikzpicture}[->,>=Stealth,thick,
  vx/.style={circle,draw=nodeblue,fill=blue!6,minimum size=6mm,font=\footnotesize},
  spam/.style={circle,draw=trapred,fill=red!7,minimum size=6mm,font=\footnotesize}]
  \node[vx] (T) at (0,0) {$T$};
  \node[spam] (s1) at (-3,-1.6) {};
  \node[spam] (s2) at (-1.5,-1.6) {};
  \node[spam] (s3) at (0,-1.6) {};
  \node[spam] (s4) at (1.5,-1.6) {};
  \node[spam] (s5) at (3,-1.6) {};
  \foreach \i in {1,2,3,4,5} { \draw[trapred] (s\i) -- (T); }
  \draw[trapred] (T) to[bend left=18] (s1.north east);
  \node at (0,-2.5) {\footnotesize farm of sybil pages all pointing at target $T$};
\end{tikzpicture}
\caption{A link farm: many cheap pages funnel their teleport-seeded mass into target $T$. Each spam page collects $(1-\alpha)/n$ from teleport and forwards its share to $T$.}
\label{ch36:fig-farm}
\end{figure}

\begin{intuitionbox}
Every page, however worthless, is \emph{guaranteed} a baseline score of at least $(1-\alpha)/n$ from teleportation. Create $f$ junk pages and point them all at $T$, and you have manufactured roughly $\alpha f \cdot (1-\alpha)/n$ of inflow for $T$ out of thin air. Worse, if $T$ links \emph{back} into the farm, mass recirculates and amplifies (a geometric series in $\alpha$). The defense the damping factor offers is the $(1-\alpha)$ leak: every cycle through the farm loses a $1-\alpha$ fraction to teleport, capping the amplification at $1/(1-\alpha)$.
\end{intuitionbox}

\begin{theorem}[Spam mass is bounded by the leak]
\label{ch36:thm-spam}
Let $S \subseteq V$ be a set of pages an adversary fully controls (a link farm), and suppose all teleport is uniform. The total PageRank captured by $S$, namely $\pi(S) = \sum_{u\in S}\pi_u$, satisfies
\[
  \pi(S) \;\le\; \frac{(1-\alpha)\,\abs{S}/n \;+\; \alpha\,\mathrm{in}(S)}{1-\alpha},
\]
where $\mathrm{in}(S) = \sum_{u\notin S}\pi_u\!\!\sum_{w\in S,\,u\to w}\!\! 1/\deg^+(u)$ is the legitimate inflow into $S$ from outside. With \emph{no} legitimate inflow ($\mathrm{in}(S)=0$), $\pi(S) \le \tfrac{\abs{S}}{n}$: a farm of $f$ pages cannot beat the score $f$ \emph{average} pages would get anyway.
\end{theorem}

\begin{proof}[Proof sketch]
Sum the stationary equation over $u \in S$. Mass entering $S$ each step is: teleport $(1-\alpha)\abs{S}/n$, plus $\alpha$ times link-mass from outside ($\alpha\,\mathrm{in}(S)$), plus $\alpha$ times link-mass that stays inside $S$ ($\le \alpha\,\pi(S)$, since at most all of $S$'s mass recirculates). At stationarity the mass on $S$ equals the mass entering it, so $\pi(S) \le (1-\alpha)\abs{S}/n + \alpha\,\mathrm{in}(S) + \alpha\,\pi(S)$; solving for $\pi(S)$ gives the claim, and $\mathrm{in}(S)=0$ yields $\pi(S)\le\abs{S}/n$.
\end{proof}

\begin{whybox}
The theorem shows how teleportation functions as an inherent defense against manipulation. Internal recirculation inside a farm is throttled by the factor $1/(1-\alpha)$; the only way to push $\pi(S)$ above the baseline $\abs{S}/n$ is to acquire \emph{legitimate} inflow $\mathrm{in}(S)$ -- links from outside pages that themselves carry authentic PageRank. Because acquiring high-reputation links cannot be automated by a spammer, manipulating the ranking becomes prohibitively difficult. A larger $\alpha$ raises the amplification cap $1/(1-\alpha)$, which is the second reason (beyond convergence speed) that production systems keep $\alpha$ moderate.
\end{whybox}

\subsection{Defenses}

\begin{itemize}
  \item \textbf{TrustRank.} Seed teleport from a hand-vetted set of trustworthy pages instead of uniformly. Trust then propagates along links; spam farms, lacking inbound links from trusted seeds, never receive the seeded mass. This is personalized PageRank with a curated $v$.
  \item \textbf{Anti-TrustRank / BadRank.} Run the chain on the \emph{reversed} graph from known-bad seeds to score how much a page links \emph{toward} spam, then demote.
  \item \textbf{Spam-mass estimation.} Compare a page's PageRank under uniform teleport against its PageRank under a trusted teleport; a large gap flags artificially inflated pages (Gyongyi--Garcia-Molina).
  \item \textbf{Topical and supervised signals.} Modern ranking folds PageRank-like graph features into learned models with hundreds of other signals, so no single graph trick dominates -- which itself dilutes the payoff of link spam.
\end{itemize}

\begin{historybox}
The arms race was real and fast. By the early 2000s ``Google bombing'' (coordinated anchor text plus links) could make a target page rank for unrelated queries; link-farm SEO was a cottage industry. TrustRank (Gyongyi, Garcia-Molina, Pedersen, 2004) and spam-mass estimation were direct responses. The lasting lesson: an open, link-derived score is inherently gameable, and robustness comes from (a) the structural $1-\alpha$ leak, (b) trust seeding, and (c) burying the signal among many others.
\end{historybox}

\section{The linear-algebra and Markov views, side by side}
\label{ch36:twoviews}

Comparing the two viewpoints side by side clarifies how linear algebra and probability inform each other:

\begin{center}
\begin{tabular}{@{}p{0.46\linewidth}p{0.46\linewidth}@{}}
\toprule
\textbf{Markov-chain view} & \textbf{Linear-algebra view}\\
\midrule
Surfer's long-run occupancy & Dominant eigenvector of $\mathbf{G}$, eigenvalue $1$\\
Ergodicity (every state reachable) & $\mathbf{G}$ strictly positive (Perron--Frobenius)\\
Mixing time of the chain & $\Oh(1/\log(1/\alpha))$ from spectral gap $1-\alpha$\\
One simulated step of the surfer & One matrix--vector product $\mathbf{G}x$\\
Discounted walk from $s$ & Neumann series $(1-\alpha)(\mathbf{I}-\alpha\mathbf{P}')^{-1}e_s$\\
Geometric$(1-\alpha)$ walk length & Eigenvalues scaled by $\alpha$\\
\bottomrule
\end{tabular}
\end{center}

\begin{aestheticbox}
Power iteration \emph{is} simulating the chain in bulk: instead of running one surfer for a long time and counting visits (a Monte Carlo estimate of $\pi$), we propagate the entire probability cloud one step at a time and let the spectral gap collapse it onto $\pi$. The same object -- the stationary distribution -- is a fixed point to the algebraist, a long-run frequency to the probabilist, and a centrality score to the network scientist. PageRank's durability comes from being correct in all three languages at once.
\end{aestheticbox}

\section{Historical Context, Controversies, and the Road Ahead}

\begin{historybox}
The theoretical foundations of PageRank began in an early twentieth-century probability dispute. In 1906, Andrey Markov introduced sequences of dependent random variables to refute Pavel Nekrasov's theological claim that the law of large numbers applied only to completely independent events. To prove his point, Markov hand-tallied twenty thousand vowels and consonants across Alexander Pushkin's verse novel \emph{Eugene Onegin}, showing that long-term stationary statistics govern systems with memory. Shortly after, Oskar Perron and Georg Frobenius established the spectral theory of non-negative matrices: any matrix with strictly positive real entries has a unique, simple maximum eigenvalue with a strictly positive eigenvector.

For nearly a century, Markov chains remained within physics and queueing theory. But in the late 1990s, the explosive growth of the World Wide Web created an unprecedented crisis in information retrieval. Early search engines ranked documents by text keyword frequency, which spammers quickly broke by stuffing popular words in invisible fonts on white backgrounds. Ranking by raw in-degree failed equally quickly, as bad actors could spin up thousands of dummy pages pointing to a scam site. At Cornell, Jon Kleinberg developed HITS to separate hubs from authorities using alternating matrix iterates. At Stanford, Larry Page and Sergey Brin realized that links are endorsements, and modeled a web surfer wandering along links at random.

Raw web graphs immediately trip up a simple random walk with two structural flaws: dangling nodes with no outgoing links leak all probability mass into oblivion, while spider traps (closed loops and dead-end cliques) swallow the walk permanently. Brin and Page's breakthrough was adding a uniform restart chance: with probability $\alpha \approx 0.85$ the surfer follows a real hyperlink, and with probability $1-\alpha$ they teleport to an arbitrary page uniformly at random. This teleportation makes the entire transition matrix strictly positive. By the Perron--Frobenius theorem, this ensures a unique, positive stationary distribution and creates a guaranteed spectral gap of $1-\alpha \approx 0.15$, allowing simple power iteration to converge to a robust global ranking in about fifty matrix multiplies regardless of web scale.
\end{historybox}

\subsection{What We Can Now Do vs.\ Honest Downsides and the Road Ahead}
\begin{connectionbox}
\textbf{Capabilities Unlocked}:
\begin{itemize}
  \item We can compute a globally consistent, spam-resistant centrality score across billions of directed vertices in near-linear time $O(m \log(1/\eps))$ using power iteration.
  \item The stationary distribution $\pi$ is unique, strictly positive, and independent of initial conditions.
  \item Personalized PageRank and localized push algorithms allow computing targeted, local similarity scores centered at a query vertex without ever materializing or traversing the global graph.
\end{itemize}

\textbf{Honest Downsides and Limitations}:
\begin{itemize}
  \item \emph{Vulnerability to coordinated Sybil attacks}: While resistant to trivial spam, PageRank is vulnerable to coordinated link farms. A network of colluding pages with high internal reciprocal linking and a modest amount of incoming external authority can trap and amplify PageRank, requiring counter-measures like TrustRank or anti-collusion filtering.
  \item \emph{Topic agnosticism and query blindness}: Global PageRank measures overall structural importance, not semantic relevance. Highly linked general hubs (like university or dictionary homepages) receive massive PageRank regardless of their relevance to a specialized technical query, requiring modern search systems to blend PageRank with semantic language models and click-through signals.
  \item \emph{Continuous global recomputation}: As the web evolves continuously with millions of added and broken links every second, updating global PageRank requires massive distributed matrix-vector computation (pioneering distributed frameworks like MapReduce and Pregel).
\end{itemize}

\textbf{Where We Go Next}:
\begin{itemize}
  \item In Chapter~37 (NP-Hard Graph Problems and Reductions), we cross a fundamental intellectual boundary. Spectral methods, random walks, and PageRank succeed because linear systems and eigenvector problems belong to the polynomial-time class $\mathbf{P}$. But many of the most vital graph problems---finding Hamiltonian tours, minimum vertex covers, graph coloring, and maximum cuts---are provably intractable ($\mathbf{NP}$-hard).
  \item In Chapter~38 (Approximation Algorithms on Graphs), we learn how to design algorithms with mathematically rigorous performance guarantees for these $\mathbf{NP}$-hard problems, using tools like semidefinite programming and randomized hyperplane rounding.
\end{itemize}
\end{connectionbox}

\section{Exercises}
\label{ch36:exercises}

\begin{exercise}[Trace with teleportation changed]
Recompute the PageRank of Figure~\ref{ch36:fig-web4} with $\alpha = 0.5$, uniform teleport. Do two power-iteration steps by hand and report $x^{(2)}$. Confirm every score moved toward $0.25$ relative to the $\alpha = 0.85$ answer, and explain why in one sentence.
\end{exercise}

\begin{exercise}[Add a dangling node]
Insert a page $E$ with an edge $A \to E$ and no out-edges of its own. Write the new sparse update rule including the dangling correction, then compute $\pi$ to two iterations. What fraction of total mass sits on $E$ at the fixed point, and where does $E$'s mass go each step?
\end{exercise}

\begin{exercise}[Eigenvalue check]
For the $4\times4$ matrix $\mathbf{P}$ of Example~\ref{ch36:ex-Pmatrix}, verify numerically that the subdominant eigenvalue of $\mathbf{G} = 0.85\mathbf{P} + 0.15\cdot\tfrac14\mathbf{1}\mathbf{1}^\top$ has modulus $\le 0.85$, consistent with Theorem~\ref{ch36:thm-eigen}. (You may use any linear-algebra library.)
\end{exercise}

\begin{exercise}[Prove the $\ell_1$ contraction directly]
Show that for any zero-sum vector $w$ ($\mathbf{1}^\top w = 0$), $\norm{\mathbf{G}w}_1 \le \alpha\norm{w}_1$, using the coupling argument from the Why-box (split $w$ into its positive and negative parts and track where each goes under one step of $\mathbf{G}$). Conclude Theorem~\ref{ch36:thm-converge} without diagonalizing.
\end{exercise}

\begin{exercise}[Personalized PageRank as a walk]
Prove that $\mathrm{ppr}_s(t) = (1-\alpha)\sum_{k\ge0}\alpha^k \Pr[\text{$k$-step walk from $s$ ends at $t$}]$ equals the stationary distribution of the teleport-to-$s$ chain. Then show $\mathrm{ppr}_s(t) = \mathrm{ppr}_t(s)$ does \emph{not} generally hold on directed graphs, but does relate via a degree factor on undirected graphs.
\end{exercise}

\begin{exercise}[Push correctness]
Prove the push invariant $\mathrm{ppr}_s = p + \sum_u r_u\,\mathrm{ppr}_u$ is preserved by one iteration of Algorithm~\ref{ch36:alg-push}. Then prove the $\Oh\!\big(1/((1-\alpha)\tau)\big)$ bound on the total degree-weighted work by amortizing against settled mass.
\end{exercise}

\begin{exercise}[Spam bound]
Complete the proof of Theorem~\ref{ch36:thm-spam} with full balance equations (do not skip the internal-recirculation term). Then construct a concrete farm of $f$ pages with one legitimate inbound link of PageRank $p_0$ and compute the exact $\pi(S)$ as a function of $\alpha, f, n, p_0$.
\end{exercise}

\begin{exercise}[Break it, then fix it]
Design a small graph (at most 8 nodes) where a link farm of 3 nodes captures more than twice its fair share $3/n$ of PageRank by exploiting legitimate inflow plus internal recirculation. Then re-rank using a TrustRank teleport vector seeded on the two legitimate nodes and show the farm's captured mass collapses.
\end{exercise}

\begin{exercise}[Monte Carlo PageRank]
Instead of power iteration, estimate $\pi$ by simulating $R$ independent surfers for $T$ steps each on the Google matrix and tallying visit frequencies. Derive how large $R$ must be for an additive $\eps$ estimate of a single $\pi_u$ with confidence $1-\delta$ (Hoeffding). Compare the total work to power iteration and state when Monte Carlo wins (hint: when you only need a few coordinates).
\end{exercise}

\begin{exercise}[Periodicity without teleport]
Build a directed graph whose raw walk matrix $\mathbf{P}$ is periodic (every node has period $2$, e.g.\ a directed bipartite cycle) so raw power iteration oscillates and never converges. Show that any $\alpha < 1$ teleport restores convergence, and quantify how the oscillation amplitude decays as $\alpha^k$.
\end{exercise}
